\documentclass[11pt,letterpaper]{article}
\usepackage[T1]{fontenc}
\usepackage[utf8]{inputenc}
\usepackage{lmodern}
\usepackage[margin=1in,headheight=14pt]{geometry}
\usepackage{amsmath,amssymb,amsthm,mathtools}
\usepackage{microtype,booktabs,longtable,array,enumitem,xcolor,listings,fancyhdr,xurl}
\usepackage[colorlinks=true,linkcolor=blue!45!black,citecolor=blue!45!black,urlcolor=blue!45!black]{hyperref}
\hypersetup{pdftitle={Counter evaluation: a smaller equality interface for iPRA},pdfauthor={Research report prepared for Vladimir Reshetnikov}}
\newtheorem{theorem}{Theorem}[section]
\newtheorem{lemma}[theorem]{Lemma}
\newtheorem{proposition}[theorem]{Proposition}
\newtheorem{corollary}[theorem]{Corollary}
\theoremstyle{definition}
\newtheorem{definition}[theorem]{Definition}
\newtheorem{remark}[theorem]{Remark}
\newcommand{\ipra}{\mathsf{iPRA}}
\newcommand{\CE}{\mathsf{CE}}
\newcommand{\K}{\mathcal K_R}
\newcommand{\TR}{T_R}
\newcommand{\Rec}{\mathop{\mathrm{Rec}}}
\newcommand{\nf}{\mathop{\mathrm{nf}}}
\newcommand{\sg}{\mathop{\mathrm{sg}}}
\newcommand{\cl}[2]{\langle #1;\,#2\rangle}
\newcommand{\code}[1]{\texttt{#1}}
\setlength{\emergencystretch}{2em}
\setlist{itemsep=3pt,topsep=5pt}
\pagestyle{fancy}\fancyhf{}
\fancyhead[L]{\small Counter evaluation for iPRA}
\fancyhead[R]{\small A complete five-pattern presentation}
\fancyfoot[C]{\thepage}
\lstset{language=Python,basicstyle=\ttfamily\scriptsize,columns=fullflexible,keepspaces=true,breaklines=true,showstringspaces=false,frame=single,rulecolor=\color{black!25},keywordstyle=\color{blue!55!black},commentstyle=\color{black!55},stringstyle=\color{green!35!black}}
\title{\textbf{Counter evaluation}\\[4pt]\large A smaller equality interface for intuitionistic primitive recursive arithmetic}
\author{Research report prepared for Vladimir Reshetnikov}
\date{November 11, 2026\\[4pt]\small Repository snapshot: \texttt{831081033253cca2115196cf86bb99c77ed2aeae}}
\begin{document}
\maketitle
\begin{abstract}
On the repository's existing counter-aware, successor-seeded recursor,
we replace separate computation and equality-transport mechanisms by
reflexivity and one uniform binary rule. From $s=t$ and
$R(Ss,a;i,u.h)=r$, the rule infers $r=h(t,R(t,a;i,u.h))$.
Two applications recover root equality composition; two applications
also recover congruence in any fixed term context. Neither derivation
uses induction, numerical cases, or a logical rule.

Together with the existing Eq/$D$ clause interface, this gives a complete
five-pattern presentation bi-interpretable with full first-order iPRA.
Both interpretation domains are whole and both numerical composites
have isomorphism graph $x=y$. The paper specifies every rule, accounts
for zero-body normalization, proves the interpretations and their
composites, and supplies standalone checkers and executable proof transformations.

The complete new checker has 3004 normalized bytes, 854 significant
tokens, and 987 AST nodes. The matched computation-plus-transport
control has 3272 bytes, 929 tokens, and 1063 nodes. This is a local
mathematical and implementation improvement, not a global minimum:
a raw two-case scalar control is five bytes smaller, and the historical
516-byte compact encoding is not surpassed. The recursor, its direct
primitive-recursion compiler, clause logic, and canonical zero convention
are attributed prior work. The new contribution is their uniform
counter-evaluation interface and its equality reconstruction.
\end{abstract}
\clearpage
\setcounter{tocdepth}{1}\tableofcontents
\clearpage

\section{What is simpler, and what is inherited}
\label{sec:comparison}

The supplied repository is unusually far along: its main index describes
fifty research packages, a Lean formalization of the native scalar basis,
and several recently submitted alternatives~\cite{index}.
The five incoming archives at the pinned snapshot were also inspected.
In particular, five inference patterns, removal of first-class $Q$,
canonical zero computation, and deletion of a root case from scalar
evaluation have already been proposed~\cite{clausal,raw,iteration}.
None is claimed here as a first discovery.

The useful new change is to put the supplied equality in the
\emph{counter} of a counter-aware recursor. Its update can project that
counter, which makes the same evaluation rule perform root equality
composition. An update that applies any fixed term template to the
counter makes it perform arbitrary congruence. Computation is the
reflexive instance of the same rule.

This differs from the existing counter-congruence presentation, which
retains separate root replacement and computation rules, and from its
wrapped-counter refinement~\cite{counter,wrapped}. It also differs from
the scalar evaluation rule, whose first premise identifies a preceding
\emph{state} rather than an arbitrary counter~\cite{eval,iteration}.

\begin{center}\small
\begin{tabular}{@{}p{.44\linewidth}rrr@{}}
\toprule
Complete standalone checker & Bytes & Tokens & AST nodes\\
\midrule
Same counter language: computation + transport & 3272 & 929 & 1063\\
\textbf{Counter evaluation (this report)} & \textbf{3004} & \textbf{854} & \textbf{987}\\
Scalar zero-normal, one-case evaluation & 3061 & 884 & 1025\\
Incoming raw scalar CE5, unchanged & 3046 & 853 & 993\\
Raw scalar CE5 with known root-case deletion & 2999 & 848 & 986\\
Incoming canonical clause checker, unchanged & 3305 & 946 & 1072\\
\bottomrule
\end{tabular}
\end{center}

The first comparison keeps the term language, normalization, formula
syntax, clause representation, validation, scoping, and public input
contract fixed. The equality interface alone changes. Its reductions
are 268 bytes (8.19\%), 75 tokens (8.07\%), and 76 AST nodes (7.15\%).
The other rows expose the actual tradeoffs rather than establishing
dominance over every existing implementation. The new checker is
smaller in all three coordinates than the scalar one-case zero-normal
control, but not than the raw two-case control. Five-pattern calculi
already existed; a smaller absolute rule count is not the claim.

All three costs count the complete checker, including normalization,
scope validation, JSON input, and conclusion checking. No producer,
translator, or project library is a trusted import. The historical
compact checker family has a different certificate interface and much
smaller source. Its 516-byte envelope remains below every readable
checker in this table~\cite{canonical}.

\section{Numerical syntax and canonical substitution}
\label{sec:syntax}

Raw numerical terms have one sort and grammar
\[
 t ::= x\mid 0\mid R(n,a;i,u.h).
\]
Only the update $h$ binds $i,u$. The first bound variable receives the
old counter, the second the preceding state; passive numerical
parameters are allowed. Define
\[
                         S(a):=R(0,a;i,u.0).
\]
The inherited intended computations are
\begin{align}
 R(0,a;i,u.h)&=Sa, \label{eq:rzero}\\
 R(Sn,a;i,u.h)&=h[n/i,R(n,a;i,u.h)/u]. \label{eq:rsucc}
\end{align}
This recursor and its elementary arithmetic compiler are prior work
in the repository~\cite{counter,scalar}. We use two bound numerical
variables in the update, unlike the later scalar iterator.

\begin{definition}[Canonical terms]
Every child of a canonical term is canonical, and an $R$ node whose
counter is literally $0$ has body literally $0$. The smart constructor is
\[
 \rho(n,a,h)=
 \begin{cases}S(a),&n\text{ is literally }0,\\
 R(n,a;i,u.h),&\text{otherwise}.
 \end{cases}
\]
Normalize a finite raw tree recursively by this constructor. Numerical
substitution means ordinary capture-avoiding substitution followed by
this normalization. All displayed formulas and rule instances use
this convention.
\end{definition}

Thus \eqref{eq:rzero} is definitional. No successor unfolding is
performed in normalization, and no test for \emph{provable} equality
with zero occurs. This is the existing canonical-zero convention,
with two update binders instead of one~\cite{canonical}.
It is part of the trusted syntax and is counted in the checker.

Normalization is idempotent and preserves scopes. It commutes with
substitution in the form
\[
 \nf(t\sigma)=\nf\bigl(\nf(t)\,\nf(\sigma)\bigr).
\]
For the $R$ case, either the final counter is zero, when both sides erase
the body, or it is not, when the structural induction hypotheses apply
to all retained children. A body erased at an earlier stage cannot
become relevant later: substitution preserves literal zero. The same
argument gives composition of normalized substitutions and extends to
the two formula constructors below. Semantic correctness under update
bodies is treated separately in Section~\ref{sec:interpretations}; it
does not follow merely from syntactic idempotence.

A conventional first-order signature may replace every fixed finite
body description by its corresponding function symbol, with its passive
parameters as arguments. This is an externally generated signature.
There is no function variable and no internal total evaluator for
arbitrary primitive-recursive program codes.

\section{The complete five-pattern calculus}
\label{sec:rules}

Object formulas and retained clauses have the forms
\[
 A::=(s=t)\mid D_{x,y}A,\qquad
 \cl{A}{y.B(y)},
\]
where the clause variable $y$ is absent from its antecedent $A$.
Contexts are finite reusable lists of clauses. The ordinary readings are
\[
 D_{x,y}B:=\exists x\bigl(B(x,0)\to\forall y\,B(x,y)\bigr),
 \qquad
 \cl{A}{y.B}:=A\to\forall y\,B.
\]
Let $H$ be $0=0$. An ordinary local assumption $A$ is represented by
the constant clause $\cl{H}{v.A}$, with $v$ fresh and unused. Its use
requires the reflexivity proof of $H$; there is no implicit hypothesis
instruction. Write $\Gamma,A$ as shorthand for adjoining this clause.

\subsection{Reflexivity and counter evaluation}
\[
 \frac{}{\Gamma\vdash a=a}\ (\mathrm{Refl}),
 \qquad
 \frac{\Gamma\vdash s=t\qquad
       \Gamma\vdash R(Ss,a;i,u.h)=r}
      {\Gamma\vdash r=h[t/i,\rho(t,a,h)/u]}\ (\CE).
 \tag{Arithmetic}
\]
The substitutions in $h$ are simultaneous and capture avoiding.
The seed, body, and passive parameters are read from the second
premise and kept fixed. In particular, occurrences inside $a$ or $h$
that happen to resemble $s$ are not replaced. The first endpoint $s$
is any term, including literal zero. The second endpoint $t$ can
also be zero, in which case $\rho(t,a,h)=Sa$.

\subsection{Use, paired elimination, and paired introduction}
\[
 \frac{\cl{A}{y.B}\in\Gamma\qquad\Gamma\vdash A}
      {\Gamma\vdash B[t/y]}\ (\mathrm{Use}),
\]
\[
 \frac{\Gamma\vdash D_{x,y}B(x,y)\qquad
       \Gamma,\cl{B(b,0)}{y.B(b,y)}\vdash F}
      {\Gamma\vdash F}\ (\mathrm{DE}),
\]
\[
 \frac{\Gamma,B(w,\eta)\vdash B(w,Sz)}
      {\Gamma\vdash D_{x,y}B(x,y)}\ (\mathrm{DI}),
 \qquad
 \eta=\begin{cases}
 z,&B\text{ is a literal equation and }w\text{ is literal }0,\\
 0,&\text{otherwise}.
 \end{cases}
\]
In DE, $b$ is fresh for the outer context, conclusion, and passive
matrix parameters. In DI, $z$ is fresh for the outer context, witness,
and passive matrix parameters. Bound variables are renamed to avoid
capture. Use requires a well-scoped numerical target. All local clauses
are discharged on exiting their indicated branch.

These are the entire primitive proof system. There is no additional
induction, successor axiom, equality schema, structural proof node,
logical connective, quantifier, or contradiction rule. Weakening,
renaming, substitution, and cut are admissible operations proved below.
The designated contradiction is the object equation $0=S0$.

The label ``five patterns'' counts DI as one syntax-directed schema,
although it has two specified modes. Its test is deliberately explicit.
It is not a claim that the system has five interchangeable instructions
under every possible rule-counting convention.

\section{Equality and computation from one evaluation rule}
\label{sec:equality}

\begin{theorem}[Arithmetic reconstruction]\label{thm:arithmetic}
Refl and CE alone derive successor computation, symmetry, transitivity,
and crossed transport in every fixed numerical template:
\[
 \frac{\Gamma\vdash s=t\qquad\Gamma\vdash F(s)=r}
      {\Gamma\vdash r=F(t)}.
\]
The constructions use no logical inference, numerical cases, or
arithmetic induction.
\end{theorem}
\begin{proof}
For successor computation take the two reflexivities $n=n$ and
$B=B$, where $B=R(Sn,a;i,u.h)$. CE concludes
$B=h[n/i,\rho(n,a,h)/u]$, exactly \eqref{eq:rsucc}.

For root composition suppose $p$ proves $s=t$ and $q$ proves $s=r$.
Set
\[
                         K_s:=R(Ss,0;i,u.i).
\]
CE applied to $q$ and $K_s=K_s$ gives $K_s=r$. CE applied to $p$
and this result gives $r=t$. Denote this derived operation by
\[
                         J(p,q): (s=t),(s=r)\longmapsto(r=t).
\]
Then $J(\mathrm{Refl}(s),p)$ is symmetry of $p$. Given $p:s=t$ and
$q:t=v$, $J(q,\mathrm{Sym}(p))$ proves $s=v$.

For congruence choose the new binders fresh and put
\[
                   B_F:=R(Ss,0;i,u.F(i)),
\]
where the update ignores $u$. CE on two reflexivities gives
$B_F=F(s)$. CE with first premise $p:s=t$ and that second premise
gives $F(s)=F(t)$. Finally apply $J$ to this congruence equation
and the supplied $q:F(s)=r$. The result is $r=F(t)$.
Every term substitution in this argument is a syntactic instance of
CE. No congruence has been presumed in order to establish congruence.
\end{proof}

\begin{proposition}[Explicit proof costs]\label{prop:cost}
Count every primitive inference, including a Refl leaf, as one node;
count repeated subtrees repeatedly. Root composition has overhead
three nodes, congruence has overhead four, and crossed transport has
overhead seven. Each supplied premise tree occurs once.
\end{proposition}
\begin{proof}
$J(p,q)=\CE(p,\CE(q,\mathrm{Refl}(K_s)))$ adds two CE nodes and
one leaf. Congruence is
$\CE(p,\CE(\mathrm{Refl}(s),\mathrm{Refl}(B_F)))$, adding four
nodes. Crossed transport is $J(\mathrm{Cong}_F(p),q)$, adding
$4+3=7$. Term labels contain copies of $F$ and the endpoints, so
these are inference-node bounds, not constant byte-size or runtime bounds.
\end{proof}

\begin{theorem}[Exact equality-interface equivalence]\label{thm:interface}
Replace Refl and CE by successor computation and one-template crossed
transport, leaving the canonical language and three logical patterns
unchanged. Every sequent is derivable in one system exactly when it is
derivable in the other.
\end{theorem}
\begin{proof}
Theorem~\ref{thm:arithmetic} gives one direction by effective replacement
of each computation or transport node.

Conversely, successor computation with counter projection gives
$K_s=s$. Applying crossed transport to two copies of this equation,
with identity template, gives $s=s$. Symmetry and transitivity then
follow from the root instances of crossed transport.
For CE, equality $s=t$ gives
\[
 R(Ss,a;h)=R(St,a;h)
          =h[t/i,\rho(t,a,h)/u].
\]
The first step is congruence in the fixed counter context; the second
is successor computation. Reverse the second supplied premise
$R(Ss,a;h)=r$ and compose. This proves the CE conclusion.
The three logical rules are unchanged. Induction on finite proof trees
therefore supplies both transformations, preserving end sequents.
\end{proof}

The use of ordinary equality in the converse is legitimate: it is
derived from the reference interface before CE is simulated. The
forward reconstruction never assumes it.

\section{The logical and induction interface is complete}
\label{sec:logic}

We give the metatheory needed to connect the five patterns to ordinary
intuitionistic first-order logic. The clause and gate constructions
are inherited~\cite{choice,clausal,raw}; their instantiation here is
included to make the full theory explicit.

\subsection{Recovering both introduction choices}

Weakening and fresh-parameter renaming are already admissible by
induction on primitive derivations: neither operation changes the
literal-zero witness test. Local-assumption cut is admissible at
this stage as well. Replace each Use of the constant clause
$\cl{H}{v.A}$ by the supplied proof of $A$, weakened into its local
context and with eigenparameters renamed. This proof transformation
does not require numerical substitution or full formula transport.

Consider an auxiliary calculus with two introductions: case introduction
from $B(w,0)$ to $B(w,Sz)$ for every matrix, and progress introduction
from $E(w,z)$ to $E(w,Sz)$ for literal equations only. The auxiliary
calculus is stable under capture-avoiding numerical substitution.

The only case-introduction instance missing from primitive DI is an
equation at witness zero. Let
\[
                         Z:=R(S0,0;i,u.0).
\]
This is not literally zero, and successor computation proves $Z=0$.
Under $E(Z,0)$, equation transport gives $E(0,0)$. Apply the supplied
case premise and transport $E(0,Sz)$ back to $E(Z,Sz)$.
Primitive DI at witness $Z$ now selects case mode and has exactly the
desired conclusion $D_{x,y}E(x,y)$. Only equation transport was used.

For progress at any witness $w$, first introduce
$D_{v,y}E(w,y)$ with unused $v$ and witness zero. This is primitive
progress mode. DE opens $\cl{E(w,0)}{y.E(w,y)}$. Under $E(w,0)$,
Use supplies $E(w,Sz)$; the now-available case introduction at witness
$w$ gives the desired $D_{x,y}E(x,y)$. Closing DE completes the proof.

Thus the two calculi have identical consequence relations. Structural
weakening, exchange, contraction, renaming, and numerical substitution
follow by proof induction in the auxiliary calculus. For clarity, a
general clause-replacement lemma has the premise
$\Delta,A\vdash B(y)$, with $y$ fresh for $\Delta,A$; every Use of
$\cl{A}{y.B}$ at $t$ is replaced by that derivation instantiated at
$t$ and local-assumption cut. These are derived operations.
The single primitive witness test itself is \emph{not} invariant under
substitution: a variable witness can become zero. Admissibility is
obtained through the auxiliary calculus, not by an invalid direct
rule-for-rule argument.

\subsection{Cases, induction, and formula transport}

From $\Gamma\vdash A(0)$ and a fresh-parameter derivation
$\Gamma\vdash A(Sz)$, case DI with an unused existential coordinate,
followed by DE and Use, proves $\Gamma\vdash A(t)$. This is numerical
case distinction for every object formula.
Replacing case DI by equation progress yields contextual equation induction:
\[
 \frac{\Gamma\vdash E(0)\qquad
       \Gamma,E(z)\vdash E(Sz)}
      {\Gamma\vdash E(t)}.
 \tag{Eq-ind}
\]
Only a literal equation is allowed as $E$; the retained context may
contain arbitrary formulas or clauses.

Full object-formula transport follows by external structural induction
on a formula template. Equation templates follow from numerical
congruence and root composition. For
$A(v)=D_{x,y}B(v,x,y)$, eliminate $A(s)$ at fresh $b$. Introduce
$A(t)$ by case DI at $b$. Transport its local base $B(t,b,0)$
backwards to $B(s,b,0)$ using the smaller-template induction hypothesis,
Use the source clause at $Sz$, and transport the result forwards to
$B(t,b,Sz)$. Close DI and DE. This is induction on finite formula
syntax in the metatheory, not induction on a compound arithmetical formula.

\subsection{Explicit logical macros}

Section~\ref{sec:arithmetic-bridge} supplies a numerical selector
$C$ satisfying $C(0,a,b)=a$ and $C(Sn,a,b)=b$ by computation alone.
For a formula $A$, define a gate $G_v A$ on equation atoms by
\[
             G_v(s=t):=(C(v,0,s)=C(v,0,t)),
\]
and commute it through $D$, freshening the binders. Then $G_0A$ is
provable and $G_{Sz}A$ is interderivable with $A$.

Define the padding $L(A):=D_{x,y}G_yA$, with $x,y$ fresh. It is
interderivable with $A$: introduce it by cases from $A$; conversely
eliminate it and Use the provable zero gate at target $1$.
Every formula is a finite chain of $D$ binders ending in an equation.
Pad the shorter of $A,B$ until their depths match. Rename corresponding
binders identically and at their equation leaves define
\[
 M_v((s=t),(p=q)):=
       (C(v,s,p)=C(v,t,q)).
\]
Carry matching $D$ binders outside. Consequently
$M_0(A,B)$ is interderivable with $A$, and $M_{Sz}(A,B)$ with $B$.

Write $Q_yA(y)$ as the object abbreviation $D_{x,y}A(y)$ with $x$
unused. Choose a fixed term $\pi$ with $\pi(1)=1$ and $\pi(2)=0$;
for example ordinary recursion starting at zero and updating
$u\mapsto C(u,1,0)$ has these two computation instances. Define
\begin{align*}
 \bot_*&:=(0=S0),\\
 A\to_* B&:=Q_yM_y(A,B),\\
 A\wedge_*B&:=Q_yG_yM_{\pi(y)}(A,B),\\
 \exists_*x\,A(x)&:=D_{x,y}G_yA(x),\\
 A\vee_*B&:=\exists_*v\,M_v(A,B),\\
 \forall_*x\,A(x)&:=A(0)\wedge_*Q_xA(x).
\end{align*}
All helper variables are fresh. These are finite definitions in Eq/$D$
syntax; no extra object-formula constructor is added.

\paragraph{Verification of the intuitionistic rules.}
Implication introduction is case DI using the derivation from $A$ to
$B$. Elimination Uses the retained conditional closure at a successor.
Conjunction introduction proves both alternatives of its selected
matrix by numerical cases; elimination Uses targets $1$ and $2$ and
the displayed computations of $\pi$. No range theorem for $\pi$ is
required: the general selector argument uses numerical cases on its value.
Existential introduction is case DI at its witness, and elimination
opens DE and Uses the zero gate at $1$ to recover the witness formula.
Disjunction introduction chooses numerical tag $0$ or $1$; elimination
eliminates the existential tag and uses numerical cases and the two
given branches. Universal introduction pairs the base instance with
the conditional closure obtained from a fresh-parameter proof;
elimination uses that base in the closure.

Finally, under $0=S0$, transport through $C(v,s,t)$ and its two
computation equations yields every equation $s=t$. External induction
on the number of $D$ binders, using case introduction, yields every
object formula. This verifies ex falso. These are precisely the usual
intuitionistic introduction and elimination rules, with their usual
discharge and eigenvariable restrictions. No excluded middle for
quantified formulas is introduced.

Numerical coverage follows by cases on the macro formula
$n=0\vee_*\exists_*m\,(n=Sm)$. Its zero and successor instances
use the displayed witnesses. The designated contradiction supplies
the separation formula $0=S0\to\bot$ in the ordinary realization.

\subsection{Ordinary realization and both formula composites}

Let $\TR$ be ordinary intuitionistic first-order logic with actual
equality over the recursor terms, the two computations (zero being
canonical), coverage, separation $0=S0\to\bot$, and contextual
literal-equation induction. Read $D$ and clauses as specified in
Section~\ref{sec:rules}. Refl, CE, Use and DE are sound by equality
and intuitionistic logic. Case DI is sound by coverage; progress DI
is sound by equation induction under the retained base assumption.

Conversely, the preceding constructions realize every rule and axiom
of $\TR$. The resulting translations are faithful in both directions.
The ordinary reading of each logical macro is provably equivalent to
the connective or quantifier it represents, by the same selector and
gate arguments. For the other composite, externally induct on Eq/$D$
syntax. In the $D$ case, open its clause at a fresh witness and build
the existential, implication and universal macros; conversely eliminate
the existential macro, recover its conditional closure, and introduce
$D$ by cases. This proves native interderivability with its round trip,
including formulas under binders and arbitrary retained assumptions.

Hence this is a reconstruction of the full first-order consequence
relation, rather than only a sound translation or a coincidence of
closed equational theorems.

\section{The inherited arithmetic bridge in a short form}
\label{sec:arithmetic-bridge}

The following terms and compiler are inherited from the counter-aware
recursor proposal~\cite{counter,scalar}. Their short proofs are a useful
reason to use this language for counter evaluation.

\begin{lemma}[Predecessor, selector, and ordinary recursion]
Define
\begin{align*}
 P(n)&:=R(n,0;i,u.i),\\
 C(n,a,b)&:=P(R(n,a;i,u.Sb)),\\
 \Rec(n,a;i,u.h)&:=P\bigl(R(n,a;i,v.S(h(i,P(v))))\bigr).
\end{align*}
Then the theory derives
\begin{align*}
 P(0)&=1,&P(Sn)&=n,\\
 C(0,a,b)&=a,&C(Sn,a,b)&=b,\\
 \Rec(0,a;h)&=a,&\Rec(Sn,a;h)&=h(n,\Rec(n,a;h)).
\end{align*}
These derivations use computation and equality only.
\end{lemma}
\begin{proof}
The $P$ equations are the two computations with counter projection.
The selector computations reduce the inner $R$ to $Sa$ or $Sb$,
then apply $P(Sx)=x$. For recursion, put
$T(n)=R(n,a;i,v.S(h(i,P(v))))$. Its base is $Sa$, while
$T(Sn)=S(h(n,P(T(n))))$. Applying $P$ proves the stated equations.
\end{proof}

Note $P(0)=1$, not zero. The standard predecessor is instead
\[
                         p(n):=C(n,0,P(n)).
\]
It satisfies $p(0)=0$ and $p(Sn)=n$. Applying $P$ to $Sx=Sy$
gives successor injectivity. The sign function
$\sg(n)=C(n,0,1)$ gives $\sg(0)=0$, $\sg(Sn)=1$; hence
$Sx=0$ contradicts separation. Coverage also gives
$n=0\vee n=S(p(n))$.

Every fixed primitive-recursive description now compiles into a finite
term: zero, successor and projections have their evident translations;
composition is substitution; primitive recursion is the displayed
$\Rec$ construction. Unlike the scalar bootstrap, this compiler
needs no pairing or finite-state predecessor construction. That
advantage belongs to the inherited counter-aware basis, not to a
new arithmetic claim here.

\section{Why equation induction gives all quantifier-free induction}
\label{sec:reflection}

For completeness we give the constructive equality-reflection step.
This avoids assuming that agreement with standard numerical evaluation
already proves equivalence of theories on arbitrary models.

Using $\Rec$, define addition $x+y$ by recursion on $y$ with update
$S$, and truncated subtraction $d(x,y)$ by recursion on $y$ with
update $p$. Thus
\[
 x+0=x,\quad x+Sy=S(x+y),\qquad
 d(x,0)=x,\quad d(x,Sy)=p(d(x,y)).
\]
Literal-equation induction proves $0+y=y$ and
$Sx+y=S(x+y)$. Inducting on $y$ in the equation
$p(d(Sx,y))=d(x,y)$ proves
\[
                         d(Sx,Sy)=d(x,y).
\]
The base is $p(Sx)=x$; the step applies $p$ to the induction equation.
Induction on $a$ therefore gives $d(Sa,a)=1$.

\begin{lemma}[Positive subtraction identifies a remainder]
For all $x,y$,
\[
             C(\sg(d(x,y)),\ x,\ y+d(x,y))=x.
 \tag{Remainder}\label{eq:remainder}
\]
\end{lemma}
\begin{proof}
Induct on $y$ in this single literal equation, keeping $x$ fixed.
At zero both selected values are $x$, using $0+x=x$ and numerical
cases on $\sg(x)$. For the step set $r=d(x,y)$ and $r'=p(r)$.
If $r'=0$, the new equation selects $x=x$. Otherwise coverage gives
$r'=Sz$. Coverage on $r$, the predecessor equations, and separation
then give $r=Sr'$. The old equation selects $y+r=x$. The new selected
sum is
\[
                  (Sy)+r'=S(y+r')=y+Sr'=y+r=x.
\]
All branch distinctions are coverage. The sole induction hypothesis
is the displayed equation, so no quantified or disjunctive induction
has been used.
\end{proof}

If $d(Sa,b)=1$, \eqref{eq:remainder} gives $b+1=Sa$. Since
$b+1=Sb$, successor injectivity gives $a=b$. The converse follows
from $d(Sa,a)=1$. Define
\[
 K(n):=C(n,1,p(n)),\qquad
 \epsilon(a,b):=\sg\bigl(K(d(Sa,b))\bigr).
\]
Coverage and separation give $K(n)=0\leftrightarrow n=1$ and
$\sg(n)=0\leftrightarrow n=0$. Therefore
\[
 \epsilon(a,b)=0\ \leftrightarrow\ a=b,
 \qquad \epsilon(a,b)=0\ \vee\ \epsilon(a,b)=1.
\]
These are internal whole-domain proofs, not checks on numerals.

For every fixed quantifier-free formula $A$, form a Boolean term $b_A$
with zero meaning true:
\begin{align*}
 b_{s=t}&=\epsilon(s,t),&b_\bot&=1,\\
 b_{A\wedge B}&=C(b_A,b_B,1),&
 b_{A\vee B}&=C(b_A,0,b_B),\\
 b_{A\to B}&=C(b_A,b_B,0).
\end{align*}
External induction on the finite formula, using the two Boolean cases,
proves $b_A=0\leftrightarrow A$ and $b_A=0\vee b_A=1$.
In the implication case this uses the already established decidability
of $A$, not excluded middle for arbitrary first-order formulas.
An induction instance for $A(n)$ is now obtained from Eq-ind for
$b_{A(n)}=0$, under its unchanged retained context. This supplies
exactly the quantifier-free induction of full first-order iPRA.

\section{The two interpretations and their round trips}
\label{sec:interpretations}

Here iPRA means intuitionistic first-order logic with equality,
one function symbol for each fixed primitive-recursive description,
its defining equations, successor injectivity and separation, and
quantifier-free induction. The numerical domain of each interpretation
below is the entire source domain. No quotient equality or proper cut
is introduced.

\subsection{From the recursor theory to iPRA}

Translate the finite syntax of a fixed body $h(i,u,\vec p)$ into a
primitive-recursive function $h^\alpha$. Introduce the corresponding
fixed primitive-recursive description
\[
 r_h(0,a,\vec p)=Sa,\qquad
 r_h(Sn,a,\vec p)=h^\alpha(n,r_h(n,a,\vec p),\vec p).
\]
Translate $R(n,a;h)$ by $r_h(n^\alpha,a^\alpha,\vec p)$, and
translate canonical zero nodes directly by successor. The translation
$\alpha:\TR\to\ipra$ leaves variables, actual equality, and ordinary
logical constructors unchanged. Computations hold by definition;
coverage is a quantifier-free induction consequence (equivalently,
use the standard primitive-recursive predecessor); separation is
an iPRA axiom; equation induction is an allowed induction instance.

\subsection{From iPRA to the recursor theory}

Translate every fixed primitive-recursive description by the compiler
in Section~\ref{sec:arithmetic-bridge}. This defines
$\beta:\ipra\to\TR$ on the same domain with actual equality.
All function defining equations hold by that compiler; successor
properties were proved using $P$ and $\sg$; quantifier-free induction
is Section~\ref{sec:reflection}. Extend identically through ordinary
first-order logic. Thus both directions interpret the \emph{full}
first-order theories.

\subsection{Body extensionality and substitution coherence}

Suppose two fixed bodies are pointwise equal under a retained context,
and their seeds are equal. Equation induction on a fresh counter
proves the equality of the corresponding recursors. The base is
successor congruence on the seeds. At the step, the two successor
computations expose the bodies; replace the previous state using
the induction equation, then use the pointwise body equality at
the counter and that state. All passive parameters remain fixed.

This lemma justifies semantic substitution when two differently
written body descriptions produce different function symbols. Literal
identity of those symbols is not assumed. It also proves correctness
of recursive zero normalization inside arbitrary enclosing recursors:
induct externally on syntax, use body extensionality for the body
child, and use the zero equation when a normalized counter is zero.
The target theory already has equation induction at this stage;
there is no circular use of a desired transport theorem.

\subsection{Both composites are internally the identity}

For a fixed $R$ term, use external induction on its finite syntax.
The induction hypotheses identify the children and the translated
body. The original $R$ trajectory and the compiled trajectory of
$r_h$ share initial value $Sa$ and update $h(n,u)$.
One literal-equation induction on a fresh counter proves their
equality at every counter, under all passive parameters. Substitution
at the original counter gives
\[
                         \TR\vdash t^{\alpha\beta}=t.
\]
For a fixed primitive-recursive description, induct externally on
that description. Zero, successor, projection and composition are
immediate. In a recursion case the original and translated functions
have equal bases and steps by the description induction hypotheses;
one iPRA equation induction proves equality everywhere. Hence
\[
                         \ipra\vdash t^{\beta\alpha}=t.
\]
The superscripts here indicate successive application in written
order. Extending externally over arbitrary first-order formulas gives
provable equivalence with each round trip. The native/ordinary
formula composites were established in Section~\ref{sec:logic}.

\begin{theorem}[Main result]
The five-pattern calculus $\K$ is bi-interpretable and equiconsistent
with full first-order iPRA. Both numerical interpretation domains are
$x=x$, equality is actual equality, and the graph of each internally
definable numerical isomorphism is $x=y$. These translations preserve
the full consequence relation, including open sequents and formulas
with unbounded quantifiers.
\end{theorem}
\begin{proof}
Compose the faithful native/ordinary translations with $\alpha$ and
$\beta$. Their formula composites are provably equivalent to the
original formulas and their numerical composites are the displayed
termwise identities. The graph $x=y$ is therefore an internal total
bijection preserving all interpreted operations and relations.
The native contradiction is $0=S0$, whose ordinary reading implies
and is implied by falsity using separation and explosion. Inconsistency
thus transfers in both directions, proving equiconsistency as a
consequence of the stronger bi-interpretability theorem.
\end{proof}

\section{Executable artifacts and verification boundary}
\label{sec:implementation}

\subsection{The public certificate contract}

The standalone checker reads one JSON object with exactly the keys
\code{goal} and \code{proof}. The goal is closed. Exit status zero
and the line \code{ACCEPT} constitute acceptance; every exception
causes rejection. No Python expression in an input is evaluated.
As with the comparison programs, a concrete Python run has recursion
and memory limits and can reject an oversized valid certificate.
Completeness is a statement about the finite formal calculus and
the checker under unbounded resources, not a promise that a fixed
machine can process every finite proof.

Terms use the integer $0$, negative free-variable levels, positive
bound indices, and three-element arrays $[n,a,h]$. Inside an $R$
body, index $1$ is the state and index $2$ the counter. A $D$ body
also binds two variables, with $y$ at index $1$ and $x$ at index $2$.
Equations are $[s,t]$ and paired formulas are $[B]$. Tree arities
therefore distinguish every syntactic constructor. Booleans are
rejected as integer variables; all scopes and canonical bodies are
validated.

An opening traversal replaces one bound slot and removes it. Since
all external variables are negative levels, opening a body with an
externally scoped term introduces no free positive index to capture.
The traversal tracks the two local binders of every $R$ or $D$ body
and reapplies the canonical constructor. CE opens the state slot and
then the counter slot, implementing simultaneous substitution.
This order and the normalization of $\rho(t,a,h)$ are essential.

The proof instructions are \code{r} (reflexivity), \code{v}
(counter evaluation), \code{u} (clause use), \code{e} (DE), and
\code{i} (DI). The latter two introduce one fresh negative level
and check their branch conclusion. A freshness escape is rejected by
the scope and branch checks. The full source is Appendix~\ref{app:checker}.

\subsection{Effective proof conversion}

The archive includes the computation-plus-transport control on the
same canonical terms and an untrusted translator to counter-evaluation
proofs. A computation node becomes CE on two reflexivities. A transport
node becomes the congruence-plus-root construction of
Proposition~\ref{prop:cost}. Use, DE and DI copy their data and translate
their premise trees recursively. End formulas are preserved literally.
This is a uniform construction on finite certificates, not a search
through semantic truth values. The output is independently checked
by the standalone target checker.

\subsection{What was checked}

The independent validation suite exercises symbolic root composition,
successor computation, arbitrary-template congruence and crossed
transport, including templates with nested recursors and passive
parameters. It checks the exact $+3$, $+4$, and $+7$ node overheads.
It also exercises logical branches, an induction proof that uses its
local induction hypothesis, and rejection after that use is removed.
Negative cases include a false equality, mismatched CE counters,
noncanonical zero bodies, Boolean and out-of-scope variables,
malformed formulas, and eigenvariable or witness escapes.

Finite numerical checks separately test the recursor, selector,
predecessor and primitive-recursion compiler. They are a supplement
to the symbolic proof arguments, not a substitute for them. The
reproduction script writes the exact test counts and machine-readable
results into \code{evidence/}. The recorded run includes 37 positive
and 15 negative public certificates, each run normally and with
Python optimization (104 process checks); 44 additional open arithmetic
and substitution checks; 29 nested-template congruence/transport pairs;
and 570 numerical/compiler checks, all with the expected outcomes.
Both proof translations additionally passed 37 closed and four open
round trips, with all 15 invalid source certificates rejected. The
comparison-control suite checks the existing root-case deletion and
the scalar zero-normal variant separately.

The new equivalence and bi-interpretation proofs are written
mathematics. The included checkers and finite tests are not a new
Lean or Rocq verification. The repository's existing Lean results
certify its own native scalar calculus; they do not automatically
certify this new primitive rule or its Python implementation.

\subsection{Reproducible source metrics}

Metrics use CPython 3.12.14, remove the module docstring, apply
\code{ast.unparse}, and append one newline. Significant tokens
exclude comments, formatting tokens, and the end marker; AST nodes
are counted by \code{ast.walk}. Gzip-9 with timestamp zero is
recorded only as supplementary data. The exact source programs and
the measurement script are included. The source-size comparisons in
Section~\ref{sec:comparison} do not omit validation, parsing, or
error handling.

\section{Further research questions}

\begin{enumerate}
\item Can the counter-evaluation interface be deeply embedded in the
repository's Lean development, with a checked equivalence to its
existing native calculus? The arithmetic reconstruction is small;
the main added work is the two-binder syntax and canonicalization.
\item Can one obtain a uniform raw-syntax counter rule with comparable
checker size, accounting explicitly for the zero computation rather
than putting it into canonical substitution?
\item Can the two counter-body instantiations be represented with a
simpler complete scope discipline, improving the transparent checker
without moving validation or normalization outside the trusted code?
\item Can an explicit proof-complexity comparison between counter CE
and scalar state evaluation quantify certificate lengths and label
sizes as well as primitive inference nodes?
\item The incoming raw-clause report's restricted successor-rigidity
principle yields zero-equation induction but leaves a faithful
equality/predecessor bootstrap unresolved~\cite{raw}. Does that
strictly narrower arithmetic principle suffice? The present report
does not solve that gap or borrow its unproved conclusion.
\item Which robust complexity measure can compare genuinely different
proof interfaces? The current byte, token, AST, constructor, and
proof-expansion coordinates disagree. A global simplicity claim would
need a specified coding model and a corresponding lower-bound theorem.
\end{enumerate}

\appendix
\section{Novelty audit and attribution}

The audit used the pinned repository index, all five incoming archive
reports or claim files, and the full relevant specifications for
counter congruence, its restricted/wrapped refinement, computational
equality, canonical recursors, initializer equality, and the raw and
canonical clause presentations. This is a scoped comparison, not a
claim to have read every archived file or established literature-wide
priority.

\begin{center}\small
\begin{tabular}{@{}p{.36\linewidth}p{.57\linewidth}@{}}
\toprule
Ingredient & Attribution / distinction\\
\midrule
Counter-aware successor-seeded $R$, $P$, $C$, and $\Rec$ & Existing successor-recursion basis; recorded in counter-congruence and scalar-successor reports.\\
Zero-body canonicalization & Canonical-recursors proposal; two numerical binders here.\\
Eq/$D$ clauses and five-pattern count & Incoming canonical-clausal and raw-clausal proposals.\\
Witness-directed single DI & Existing witness-rule proposal.\\
Scalar root-case deletion & Incoming equality-from-iteration report; not rediscovered here.\\
Refl plus uniform counter CE & Proposed interface of this report; root composition and congruence follow from CE itself, without separate root replacement or computation axioms.\\
\bottomrule
\end{tabular}
\end{center}

The included baseline checker files are retained under the repository's
MIT No Attribution license. Mathematical attribution is retained
independently of that license. Source references and SHA-256 checksums
of the delivered artifacts are supplied for reproducibility.

\section{Complete standalone checker}\label{app:checker}
The following listing is the entire trusted implementation. The
separate producers, translators, comparisons, and tests are not imported.
\begin{lstlisting}
"""Counter-aware canonical recursor with one evaluation scheme.

Terms are 0, negative free levels, positive bound indices, and [n,a,h].
A recursor body binds state u at index 1 and counter i at index 2.
R(0,a;h) is canonical only with h=0; successor(a) = [0,a,0].
Formulas are [s,t] (equality) and [B] (paired binder D).
Input is {'goal': a closed formula, 'proof': a finite proof tree}.
"""
import json
import sys


def fail():
    raise ValueError('invalid certificate')


def rec(n, a, h):
    return [n, a, 0 if n == 0 else h]


def successor(a):
    return [0, a, 0]


def term(a, scope, depth=0):
    if type(a) is int:
        return -scope <= a <= depth
    return (type(a) is list and len(a) == 3
            and (a[0] != 0 or a[2] == 0)
            and all(term(t, scope, depth + 2 * (i == 2))
                    for i, t in enumerate(a)))


def formula(a, scope, depth=0):
    return (type(a) is list
            and ((len(a) == 2 and all(term(t, scope, depth) for t in a))
                 or (len(a) == 1 and formula(a[0], scope, depth + 2))))


def opening(a, value, depth=0):
    if type(a) is int:
        return value if a == depth + 1 else a - (a > depth + 1)
    size = len(a)
    result = [opening(t, value, depth + 2 * (size != 2 and i == size - 1))
              for i, t in enumerate(a)]
    return rec(*result) if size == 3 else result


def derive(proof, context, scope):
    if type(proof) is not list or not proof or type(proof[0]) is not str:
        fail()
    tag = proof[0]
    if tag == 'u' and len(proof) == 4:
        index, witness, premise = proof[1:]
        if (type(index) is not int or not 0 <= index < len(context)
                or not term(witness, scope)):
            fail()
        antecedent, body = context[index]
        if derive(premise, context, scope) != antecedent:
            fail()
        return opening(body, witness)
    if tag == 'r' and len(proof) == 2:
        a = proof[1]
        if not term(a, scope):
            fail()
        return [a, a]
    if tag == 'v' and len(proof) == 3:
        s, t = derive(proof[1], context, scope)
        b, r = derive(proof[2], context, scope)
        c, a, h = b
        if c != successor(s):
            fail()
        return [r, opening(opening(h, rec(t, a, h)), t)]
    if tag in ('e', 'i') and len(proof) == 4:
        goal, data, branch = proof[1:]
        if not formula(goal, scope):
            fail()
        fresh = -scope - 1
        if tag == 'e':
            paired = derive(data, context, scope)
            if len(paired) != 1:
                fail()
            body = opening(paired[0], fresh, 1)
            clause = (opening(body, 0), body)
            target = goal
        else:
            if len(goal) != 1 or not term(data, scope):
                fail()
            body = opening(goal[0], data, 1)
            eta = fresh if len(goal[0]) == 2 and data == 0 else 0
            clause = ([0, 0], opening(body, eta))
            target = opening(body, successor(fresh))
        if derive(branch, context + [clause], scope + 1) != target:
            fail()
        return goal
    fail()


def verify(data):
    if type(data) is not dict or set(data) != {'goal', 'proof'}:
        fail()
    if not formula(data['goal'], 0):
        fail()
    if derive(data['proof'], [], 0) != data['goal']:
        fail()
    return True


if __name__ == '__main__':
    try:
        verify(json.load(sys.stdin))
    except Exception:
        print('REJECT')
        sys.exit(1)
    print('ACCEPT')
\end{lstlisting}

\clearpage
\begin{thebibliography}{99}\small\raggedright
\bibitem{index} ProveIt contributors.
\emph{Concise presentations of intuitionistic primitive recursive arithmetic}.
README at commit \texttt{831081033253cca2115196cf86bb99c77ed2aeae}.
All repository paths below are relative to
\path{Logic/PrimitiveRecursiveArithmetic/ConciseIPRA/} at this snapshot.
\url{https://github.com/VladimirReshetnikov/ProveIt/tree/831081033253cca2115196cf86bb99c77ed2aeae/Logic/PrimitiveRecursiveArithmetic/ConciseIPRA}.
\bibitem{counter} ProveIt contributors.
\emph{Counter congruence recovers unrestricted equality transport}.
\path{Research/counter_congruence/theory/counter_theorem.tex}.
\bibitem{wrapped} ProveIt contributors.
\emph{A narrower congruence rule and its exact boundary}.
\path{Research/witness_rule/theory/restricted_congruence.tex}.
\bibitem{scalar} ProveIt contributors.
\emph{Scalar successor}.
\path{Research/scalar_successor/report.tex}, especially the attribution
of the earlier counter-aware successor-recursion basis.
\bibitem{eval} ProveIt contributors.
\emph{Computational equality}.
\path{Research/computational_equality/} and
\path{Lean/ConciseIPRA/NativeEval.lean}.
\bibitem{canonical} ProveIt contributors.
\emph{Canonical recursors}.
\path{Research/canonical_recursors/README.md} and
\path{theory/canonical_R.tex} within that package.
\bibitem{clausal} ProveIt contributors.
\emph{Clausal canonical presentation}.
\path{docs/incoming/ipra_clausal_canonical.zip}.
\bibitem{raw} ProveIt contributors.
\emph{Five rules for iPRA with ordinary substitution}.
\path{docs/incoming/ipra_raw_clausal_2026-11-11.zip},
especially the rule specification, finite-seed obstruction, and
restricted successor-rigidity discussion.
\bibitem{iteration} ProveIt contributors.
\emph{Equality from iteration}.
\path{docs/incoming/ipra_equality_from_iteration_2026-11-11.zip}.
\bibitem{choice} ProveIt contributors.
\emph{Complementary introductions with one local assumption} and
\emph{From literal-equation induction to full intuitionistic PRA}.
\path{Research/certified_zero/reference/inherited_choice_basis.tex}
and \path{equality_reflection.tex}.
\bibitem{seed} ProveIt contributors.
\emph{Equality in the initializer}.
\path{Research/initializer_equality/sections/seed_rule.tex}.
\end{thebibliography}
\end{document}
